\section{The Real Objection to Intelligent Design}

Intelligent Design, in the broadest sense of the term, is the methodological principle that causation may be the result of the intention of an intelligent agent. At first glance it is difficult to understand why that should be controversial. In everyday life, we use that principle all the time. For example, when we get home, we notice that the door is unlocked or the toaster oven is on, then we conclude that someone had left the door unlocked or neglected to turn off the oven.

In point of fact, this is the most obvious form of causation, since we have the direct internal experience of a will causing events to occur. Thus, early human experience brought about various theories of theism to explain things. In August Comte's view this involved an historical progression involving

\begin{enumerate}


\item \textbf{Animism} --- everything has its own will to initiate events


\item \textbf{Polytheism} --- multiple centers of will


\item \textbf{Monotheism} --- single center of will

\end{enumerate}


The general trend of modern philosophy has been the elimination of teleological explanations which reaches its culmination in Positivism. Positivism, at least in its Comtean form, eliminated all theistic explanations in the name of facts and the laws relating them to each other.Superficially, this seems an improvement, since it claims to be objective and it rejects all the superstitions and rival claims of the various theistic systems. Thus, it is accepted by self-designated "educated", "sophisticated", or "progressive" people as a given, beyond any challenge

Although positivism may work well in describing the motions of the planets or chemical processes, the real problem arise when the it is applied to social, political, or religious issues. This requires the hypothesis of nonhuman forces to explain the group activities of human agents. Thus, with Marxism, there is the process of dialectical materialism that "really" explains what's going on, rather than the illusory intentions of the human actors. More recently the "selfish gene" is used as an explanatory principle. Of course, all such theories require the "scientist" to be able to step outside the explanatory principle, as it were, so that his theory is not caught up in a vicious circle.

Such lines of reasoning have the advantage that arguments about morality, or politics, or religion are beside the point. All that is needed is to identify the hidden cause: e.g., genes, economic structures, sexual repression, and so on. This has reached its synthesis, largely under the influence of Ken Wilber, in the idea that as a person "evolves", his moral perspective will change. The claim, then, is that a man's evolution can be precisely measured by his score on a test of responses to various social and moral questions. Thus, someone with some other view is not simply mistaken, but he is much worse: an "unevolved" being. Of course, the very highest beings manage to integrate somehow all the lower levels, but don't count on in it.

The net effect is that the real meaning of human actions is disguised in a web of impersonal explanations or appeals to higher consciousness. The general form of the argument is being unfolded:

\begin{itemize}


\item Truth is betrayed by partisanship in an open society

\item Truth is replaced by unquestioned worldviews

\item Worldviews are justified by claims of metaphysical force or personal evolution
\end{itemize}


This will all come together in the forthcoming review of \textit{The Machiavellians: Defenders of Freedom} by James Burnham.


\flrightit{Posted on 2008-04-17 by Cologero}

